La integración conecta el procesador con los bloques que permiten cargar
programas, controlar su ejecución y consultar su estado desde una PC. Este
capítulo describe cómo colaboran esos bloques y qué pruebas comprueban su
funcionamiento en conjunto.

\section{Composición del sistema}

El sistema funcional está formado por cuatro bloques principales: la UART, el
controlador de protocolo, el serializador de respuestas y el procesador. La
tabla~\ref{tab:bloques-sistema} presenta la función de cada uno dentro del
recorrido de una solicitud.

El módulo \code{system_top} conecta estos bloques entre sí. A su vez,
\code{basys3_top} adapta el sistema a la placa: conecta sus pines, genera el
reloj de funcionamiento y acondiciona la señal de reset.

\begin{table}[H]
  \centering
  \small
  \begin{tabularx}{\textwidth}{L{4.4cm}Y}
    \toprule
    \textbf{Bloque} & \textbf{Función en la integración} \\
    \midrule
    \code{uart_core} & Sincroniza RX, genera el tick de baudios, recibe y
    transmite bytes y conserva las FIFO de recepción y transmisión. \\
    \addlinespace
    \code{protocol_controller} & Interpreta \code{LOAD}, \code{RUN},
    \code{STEP}, \code{RESET} y \code{CHECK_READY}; coordina las solicitudes,
    los tiempos de espera y la observación. \\
    \addlinespace
    \code{response_serializer} & Captura el descriptor de respuesta y produce
    la cabecera, el contexto, el banco de registros, los latches y la lista de
    DMEM en el orden aprobado. \\
    \addlinespace
    \code{processor_top} & Ejecuta la imagen y entrega el estado observable:
    PC, contador, RF, latches y metadatos de fault. \\
    \bottomrule
  \end{tabularx}
  \caption{Bloques de la integración de sistema.}
  \label{tab:bloques-sistema}
\end{table}

En un intercambio, la UART recibe los bytes de la PC y los entrega al
controlador. Este interpreta la solicitud y la coordina con el procesador;
después, el serializador forma la respuesta y la UART la transmite de vuelta.
El controlador atiende una operación normal por vez y, durante Debug, mantiene
estable el estado que se está enviando hasta que termina la respuesta.

\section{Coordinación de comandos y observación}

\subsection{LOAD y preparación}

La PC inicia \code{LOAD} enviando una cabecera de seis bytes que incluye el
tamaño del programa. El controlador valida ese tamaño y solicita al control
de ejecución que acepte la carga. Una vez recibida la respuesta
\code{ACCEPTED}, la PC transmite el contenido y el cargador del procesador
reconstruye las palabras para escribirlas en IMEM.

Cuando se recibe la última palabra, comienza la preparación del estado
inicial: se inicializan el PC, el banco de registros, los registros interetapa
y el contador, y se limpian las marcas de validez de DMEM. Al terminar la
preparación se publica la imagen y el sistema responde \code{DONE}, indicando
que el programa está disponible para ejecutar.

\subsection{RUN, STEP y Debug}

Una solicitud \code{RUN} aceptada inicia la ejecución continua y recibe una
respuesta \code{ACCEPTED}. Cuando el programa termina por \code{halt} o por un
fault y finaliza el drenado, la FPGA envía \code{FINISHED} o \code{FAULT} junto
con el snapshot del estado final. En cambio, \code{STEP} ejecuta un ciclo
efectivo y responde \code{DONE} con el estado posterior a ese ciclo.

Para construir el snapshot, el serializador consulta el estado que el
procesador expone para Debug y recorre el banco de registros y DMEM mediante
sus caminos de lectura compartidos. Mientras se transmite la respuesta, el
controlador mantiene retenida la CPU y excluye nuevas operaciones que podrían
modificar ese estado. Así, todos los campos enviados corresponden a la misma
observación, hasta que se transmite el bit de parada del último byte.

\section{Comprobaciones de integración}

Las pruebas de integración comprueban que los bloques conectados respeten el
protocolo y que las respuestas coincidan con el estado del procesador. Para
ello se simulan solicitudes y respuestas como tramas UART de formato 8N1 y se
comparan los bytes transmitidos con los efectos observados en el sistema.

La tabla~\ref{tab:casos-integracion} enumera los grupos del catálogo y el
tema principal de cada uno. Los dos primeros comprueban la UART que sirve de
base a los intercambios; los restantes abarcan el protocolo, la ejecución,
Debug y las condiciones de recuperación.

\begin{table}[H]
  \centering
  \small
  \begin{tabular}{llr}
    \toprule
    \textbf{Grupo} & \textbf{Tema} & \textbf{Casos} \\
    \midrule
    \code{TEST-UART-001} & Temporización UART, sincronización y tramas & 31 \\
    \code{TEST-UART-002} & FIFO de recepción y sus indicadores & 21 \\
    \code{TEST-UART-003} & Comandos, disponibilidad y rechazos & 44 \\
    \code{TEST-UART-004} & Carga de programas & 18 \\
    \code{TEST-UART-005} & Tiempo de espera y errores de carga & 24 \\
    \code{TEST-UART-006} & Recuperación por silencio y error de comando & 8 \\
    \code{TEST-UART-007} & Snapshot fijo, registros y pipeline & 7 \\
    \code{TEST-UART-008} & Lista de DMEM & 10 \\
    \code{TEST-UART-009} & TX llena y finalización de transmisión & 9 \\
    \code{TEST-UART-010} & Comandos RUN y STEP & 5 \\
    \code{TEST-UART-011} & Exclusión durante Debug & 31 \\
    \code{TEST-UART-012} & Orden de respuestas y reset & 14 \\
    \code{TEST-UART-013} & Secuencia de reset & 16 \\
    \code{TEST-UART-014} & Tiempos del enlace integrado & 10 \\
    \midrule
    \multicolumn{2}{r}{Total UART/Debug} & 248 \\
    \midrule
    \code{TEST-PLAT-001} & Perfiles y capacidades & 15 \\
    \code{TEST-PLAT-003} & Composición RTL y dependencias & 1 \\
    \bottomrule
  \end{tabular}
  \caption{Grupos de la integración de sistema.}
  \label{tab:casos-integracion}
\end{table}

Los 248 casos de UART y Debug figuran como implementados en el ejecutor. Además
de las pruebas de simulación, \code{TEST-PLAT-001} comprueba configuraciones de
las memorias y \code{TEST-PLAT-003} examina la composición del sistema y sus
dependencias. Los niveles de prueba y el significado de estos estados se
presentaron en el capítulo~\ref{chap:verificacion}.

En la regresión general se realizaron 2695 ejecuciones RTL: 2692 aprobaron en
la primera corrida y las tres restantes aprobaron al repetirse. Estas últimas
correspondieron a las variantes \code{TEST-UART-006/C001}, \code{C013} y
\code{C014}, que comprueban la recuperación de la comunicación.

\subsection{Carga, observación y reset}

Para comprobar la integración se realizaron 44 ejecuciones dirigidas de los
grupos \code{TEST-UART-004}, \code{TEST-UART-008} y \code{TEST-UART-013}, que
abarcan la carga de programas, la lista de memoria utilizada y la secuencia de
reset. Las 44 aprobaron con Icarus Verilog 13.0.

Además, \code{TEST-UART-013/C001} aprobó en XSim con \code{basys3_top}, el
generador de reloj, la CPU y la UART configurada a \num{19200} baud. Esa prueba
comprueba la respuesta al reset y el estado posterior del sistema utilizando
los componentes integrados de la plataforma.

\subsection{Revisión de la composición}

La revisión \code{TEST-PLAT-003} examinó la composición del procesador y del
sistema, incluidas sus conexiones y dependencias combinacionales. En la
implementación se comprobó que \code{system_top} conservara las instancias de
CPU, UART, controlador y serializador, junto con los componentes de reloj y
reset de la plataforma.

Los informes de Vivado no registraron puntos internos sin restricción temporal,
registros sin reloj ni lazos combinacionales o latches. El análisis de recursos
y temporización se desarrolla en el capítulo~\ref{chap:plataforma-timing}.

\section{Registro de evidencia}

Para poder revisar un resultado, cada ejecución identifica el caso, el nivel de
verificación y la configuración utilizada. Según lo que compruebe la prueba,
la evidencia puede incluir:

\begin{itemize}
  \item las solicitudes y respuestas intercambiadas por UART;
  \item el snapshot con los registros, el PC, el contador, el estado del
    pipeline, los datos de fault y la memoria utilizada;
  \item los efectos observados ante una carga parcial, un error de recepción o
    un tiempo de espera agotado;
  \item la secuencia de señales y respuestas asociada a un reset.
\end{itemize}

Estos datos permiten comparar la respuesta esperada con la observada y, en las
pruebas que incluyen Debug, comprobar que el snapshot corresponda al estado
posterior al último ciclo autorizado. Los resultados se conservan en el
\evidencefile{tp3-issue-17-evidencia.zip}{paquete de evidencia}.

\section{Resultado de la integración}

Las pruebas comprueban el recorrido de las solicitudes y respuestas entre la PC
y el procesador, así como la coordinación necesaria para que Debug observe un
estado coherente. También se ejercitan la recuperación ante errores, el reset
y la carga de nuevos programas.

La comprobación en la Basys 3 complementa las simulaciones: 150 de las 151
variantes seleccionadas aprobaron en placa. La variante
\code{TEST-UART-013/C009}, correspondiente a reset durante la preparación, quedó
pendiente de ejecución física y cuenta con un resultado aprobado en simulación
de sistema. El alcance de esa validación se presenta en el
capítulo~\ref{chap:aceptacion-fisica}.
